\documentclass[12pt,a4paper]{article}

% Biên dịch bằng XeLaTeX trên Overleaf.
\usepackage{fontspec}
\setmainfont{TeX Gyre Termes}
\usepackage[margin=2.5cm]{geometry}
\usepackage{enumitem}
\usepackage{microtype}

\setlength{\parindent}{0pt}
\setlength{\parskip}{0.35em}
\setlist[enumerate]{leftmargin=2em, itemsep=0.25em, topsep=0.35em}

\begin{document}

\textbf{Nhóm:} 1
\hfill
\textbf{Thành viên:} Tô Thành Nguyên
\hfill
\textbf{Ngày báo cáo:} 20/09/2026

\vspace{1em}

\textbf{Những nội dung đã làm trong 2 tuần qua:}

\textbf{Sprint 3: Tổ chức context và bộ nhớ phiên (07/09 -- 20/09)}

\begin{enumerate}
    \item Xây dựng context ba tầng gồm quy tắc hệ thống, trạng thái phiên và lịch sử hội thoại gần nhất trước mỗi lần gọi mô hình.

    \item Áp dụng ngân sách context mặc định 12.000 ký tự; hệ thống dừng trước khi gọi mô hình nếu phần nội dung bắt buộc vượt ngân sách.

    \item Lưu đầy đủ kết quả công cụ trong SQLite nhưng chỉ đưa excerpt JSON tối đa 1.200 ký tự cùng \texttt{tool\_call\_id} vào context.

    \item Xây dựng task card thể hiện mục tiêu, tiêu chí hoàn thành, bước hiện tại, công cụ đã dùng, kết quả và phần việc còn lại.

    \item Gắn Task, Step, ToolCall và Trace với đúng người dùng và request; chỉ tạo task khi agent thực sự gọi tool.

    \item Bổ sung reservation theo người dùng, cooperative cancellation qua \texttt{/cancel} và checkpoint trước mỗi lần gọi mô hình hoặc công cụ.

    \item Hoàn thiện migration Sprint 3, bật foreign key cho SQLite và cập nhật summary cùng watermark bằng compare-and-update nguyên tử.

    \item Xây dựng compaction theo page hữu hạn, giữ tối đa bốn exchange gần nhất và dùng extractive fallback khi mô hình tóm tắt lỗi.

    \item Hoàn thiện \texttt{/memory}, \texttt{/forget} và \texttt{/forget confirm}; memory được xóa theo ownership nhưng user và lịch sử usage vẫn được giữ.

    \item Xây dựng bộ đánh giá gồm tám tình huống trên ba chiến lược \texttt{full}, \texttt{window10} và \texttt{sprint3}; offline evaluation hoàn tất 48/48 lượt.

    \item Live evaluation schema v2 hoàn tất 24/24 lượt: \texttt{full} đạt 8/8 với 255.950 tokens, \texttt{window10} đạt 4/8 với 138.819 tokens và \texttt{sprint3} đạt 7/8 với 175.878 tokens.

    \item Chiến lược Sprint 3 giảm 80.072 tokens, tương đương 31,3\% so với \texttt{full}, và không có unsupported claim trong bộ đánh giá live.

    \item Kết quả kiểm tra cuối: Ruff đạt, 77 bài kiểm thử pytest đạt, offline artifact và live artifact đều có \texttt{incomplete=false}.

    \item Hoàn thành demo Telegram thật: nhớ dữ kiện sau hội thoại dài, gọi tool, hiển thị task card và memory, xóa phiên có xác nhận, đồng thời giữ số liệu usage.
\end{enumerate}

\textbf{Những nội dung sẽ làm trong 2 tuần tới:}

\textbf{Sprint 4: Core Toolset qua Telegram (21/09 -- 04/10)}

\begin{enumerate}
    \item Mở rộng Tool Registry với purpose, hướng dẫn sử dụng, risk level, input schema và output schema.

    \item Thêm \texttt{ToolContext} chứa identity và ownership đáng tin cậy; chuẩn hóa \texttt{ToolResult} và validate cả input lẫn output bằng Pydantic.

    \item Loại \texttt{read\_file} path-based khỏi production registry; tài liệu chỉ được truy cập qua attachment capability không để model tự truyền đường dẫn.

    \item Tích hợp Tavily Search và Extract để tìm kiếm web, lấy nguồn thật và trả citation/URL có trong tool result.

    \item Bổ sung xử lý lỗi thiếu API key, timeout, upstream reject, sai tham số và dữ liệu đầu ra không hợp lệ; không tạo kết quả giả.

    \item Xử lý PDF, DOCX và UTF-8 TXT tối đa 10 MiB qua Telegram; trích xuất theo page, section hoặc paragraph có locator.

    \item Từ chối trung thực PDF scan không có text, file hỏng, file quá lớn hoặc định dạng không hỗ trợ; chưa triển khai OCR trong Sprint 4.

    \item Xây dựng reminder theo user với các thao tác tạo, cập nhật và hủy; lưu dữ liệu UTC bền vững trong SQLite.

    \item Dùng APScheduler để phục hồi reminder sau khi bot restart, retry delivery có giới hạn và ghi trạng thái gửi rõ ràng.

    \item Kiểm thử toàn tuyến qua Telegram: tìm web, đọc nguồn hoặc tài liệu, tạo reminder, restart bot và nhận thông báo đúng user.

    \item Chạy toàn bộ pytest và Ruff sau integration; cập nhật README, \texttt{.env.example}, hướng dẫn vận hành và báo cáo Sprint 4.
\end{enumerate}

\textbf{Những trở ngại:}

\begin{enumerate}
    \item Strategy \texttt{sprint3} chưa đạt E-05 vì fact nằm giữa tool payload và không xuất hiện trong bounded head/tail excerpt.

    \item Live evaluation hiện chỉ có tám tình huống và một lần lặp nên kết quả mang tính thăm dò, chưa chứng minh tối ưu phổ quát.

    \item Web research phụ thuộc Tavily API key, quota và độ ổn định của dịch vụ bên ngoài.

    \item Parser tài liệu phải kiểm soát file nén, dung lượng mở rộng, thời gian chạy và bộ nhớ để tránh file xấu làm treo worker.

    \item Telegram không cung cấp idempotency cho việc gửi reminder; Sprint 4 dùng at-least-once delivery nên vẫn phải công khai rủi ro gửi trùng khi tiến trình crash đúng thời điểm.
\end{enumerate}

\end{document}
